\section{ConvexTok: certified path vocabularies}
\label{convextok:section}
The words ``token'' and ``graph'' have two different roles here. A TokenGT
position represents a node or edge with features and incidence metadata.
A tokenisation position instead represents a boundary in a word; an edge
jumps across a contiguous substring. ConvexTok, introduced by Tempus,
Whittington, Schmidt, Komm and Pimentel,
\emph{Tokenisation via Convex Relaxations},
\href{https://arxiv.org/html/2605.22821v1}{arXiv:2605.22821v1},
optimises the second representation. We derive its small finite
formulation, reproduce its six-word example, and give new certified
examples for finite arrow-word corpora of our infinite surface algebras.
This supplies an auditable compression layer; it does not identify the two
kinds of token or change a trained graph model.

\subsection{A vocabulary controls many occurrences}
Let $\mathcal D=(w_1,\ldots,w_N)$ be a finite multiset of nonempty words over
a mandatory alphabet $\Sigma$. For each word introduce boundary vertices
$0,\ldots,|w_n|$. Include an edge $(i,j)$ for each allowed contiguous token
$w_n[i:j]$. All single-symbol edges are free in the vocabulary budget,
but cost one each in the segmentation objective. Longer candidates form
a finite set $C$. Equal strings have the same colour even when they occur
in different words or positions. Candidates never cross word boundaries.

Write $x_e$ for occurrence flow and $c_t$ for selection of token type $t$.
Use incidence entries $-1$ at an edge's source and $+1$ at its target.
For each word let $d$ be $-1$ at its initial boundary and $+1$ at its final
boundary. Separate unit-flow equations are equivalent to joining adjacent
word terminals as in the paper: each word boundary must carry one unit,
and there are no edges across words. The finite programme is
\begin{equation}\label{convextok:lp}
 \begin{aligned}
 L_K=\min\quad &\sum_e x_e\\
 \text{subject to}\quad &Bx=d,\qquad 0\le x_e,c_t\le1,\\
 &x_e\le c_t\quad(e\text{ an occurrence of }t\in C),\qquad
 \sum_{t\in C}c_t\le K.
 \end{aligned}
\end{equation}
Requiring $x_e,c_t\in\{0,1\}$ gives the integer problem. On the acyclic
position graph its flow is one path per word, hence one segmentation.
Conversely any segmentation with at most $K$ extra token types gives such
an integer point. Consequently the LP optimum is a lower bound for the
integer optimum $I_K$. The constraint is per occurrence: selecting
$\mathtt{ab}$ once permits both occurrences in $\mathtt{abab}$, whose
cost is two. Replacing these constraints by
$\sum_{e:\,t(e)=t}x_e\le c_t$ would incorrectly prohibit that reuse.

For a fixed discrete vocabulary $V\supseteq\Sigma$, the exact segmenter is
the backward recurrence
\[
 D(n)=0,\qquad
 D(i)=1+\min_{t\in V:\,w[i:i+|t|]=t}D(i+|t|).
\]
It returns a shortest path and a decoding by concatenation. Longest-first
greedy need not do so: for $w=\mathtt{abcd}$ and extra vocabulary
$\{\mathtt{ab},\mathtt{bcd}\}$ it returns
$\mathtt{ab}|\mathtt c|\mathtt d$, whereas the optimum is
$\mathtt a|\mathtt{bcd}$. This segmentation issue is separate from BPE's
greedy construction of the vocabulary itself.

\subsection{An exact lower bound needs a feasible certificate}
A floating solver's proposed objective alone is insufficient for an
exact lower-bound claim. Our original CPU checker writes the LP in the
form
\[
 \min q^{\mathsf T}v,\qquad Av=b,\quad Gv\le h,\quad v\ge0,
\]
where $v=(x,c)$ and $G$ includes every upper-bound row $v_i\le1$.
It uses SciPy/HiGHS from the already approved environment to propose
primal and dual solutions, reconstructs rational entries, and then checks
\emph{every} condition with exact rational arithmetic:
\[
 Av=b,\quad Gv\le h,\quad v\ge0,\qquad
 z\le0,\quad A^{\mathsf T}y+G^{\mathsf T}z\le q,
 \qquad q^{\mathsf T}v=b^{\mathsf T}y+h^{\mathsf T}z.
\]
These are executable rational witnesses, not a claim that floating-point
optimisation itself is a proof assistant.
\begin{proposition}[Finite rational certificate]
Every feasible primal $v'$ obeys
$q^{\mathsf T}v'\ge b^{\mathsf T}y+h^{\mathsf T}z$ for any dual pair
satisfying the displayed inequalities. Equality with one feasible $v$
certifies the exact LP optimum.
\end{proposition}
\begin{proof}
Since $v'\ge0$, the componentwise dual inequality gives
\[
 q^{\mathsf T}v'\ge y^{\mathsf T}Av'+z^{\mathsf T}Gv'
 \ge y^{\mathsf T}b+z^{\mathsf T}h.
\]
The second inequality uses $z\le0$ and $Gv'\le h$. Equality therefore
matches a lower and an upper bound.
\end{proof}
Independently, our checker exhausts every size-$K$ subset of $C$ and runs
the shortest-path segmenter on each. This covers the at-most-$K$ problem,
because adding candidate tokens cannot increase a shortest-path cost.
For $K>|C|$, all candidates suffice. Every stored witness is checked;
deliberately perturbed primal and dual witnesses are rejected.

\subsection{Three rounding rules and worked numerical gaps}
The paper's Det rule chooses the $K$ largest $c_t$; Bias chooses the
$K$ largest $c_t/|t|$; Int retains only values essentially equal to one.
Our certificates are rational, so Int uses exact equality to one and
Det/Bias use lexicographic tie-breaking. After choosing tokens we
recompute a discrete shortest path; a fractional flow is not itself a
tokenisation. An integral colour in one LP solution need not be present
in every optimum. For example, with word $\mathtt{abc}$ and candidates
$\mathtt{ab},\mathtt{bc}$, either single-token choice is optimal.

The paper's example is
\[
 \mathcal D_0=(\mathtt{abc},\mathtt{abd},\mathtt{abe},
              \mathtt{bc},\mathtt{bd},\mathtt{be}),\qquad K=3.
\]
The mandatory-symbol count is fifteen. BPE first creates $\mathtt{ab}$,
saving three positions, then two merges save only one each: its cost is
ten. Choosing $\mathtt{bc},\mathtt{bd},\mathtt{be}$ instead gives nine.
Our exact primal/dual pair has value nine, and all $\binom73=35$
candidate vocabularies confirm $L_3=I_3=9$. All three rounding rules give
nine on the saved LP solution.

Our first new example is
\[
 \mathcal D_1=(\mathtt{cbcbc},\mathtt{cbca},
              \mathtt{cbabac},\mathtt{caaab}),\qquad K=2.
\]
One optimal fractional vocabulary has
\[
 c_{\mathtt{cbc}}=\tfrac12,\qquad
 c_{\mathtt{caaab}}=\tfrac12,\qquad
 c_{\mathtt{cbabac}}=1,
\]
with all other values zero. Its per-word costs are $3,3,1,3$, totalling
ten; the rational dual certifies optimality. The integer optimum is
eleven, witnessed by vocabulary
$\{\mathtt{caaab},\mathtt{cbabac}\}$ with costs $5,4,1,1$ and confirmed
by all $\binom{30}{2}=435$ vocabularies. Det and Bias also yield eleven,
while Int yields fifteen. The LP is therefore genuinely fractional even
though its objective is an integer.

A second original corpus,
\[
 (\mathtt{aaaab},\mathtt{baaccb},\mathtt{aac},
   \mathtt{babaab},\mathtt{cbbbca}),\qquad K=2,
\]
has $L_2=31/2<I_2=16$. Its nonzero colour weights are
$c_{\mathtt{aa}}=c_{\mathtt{babaab}}=1/2$ and
$c_{\mathtt{cbbbca}}=1$. The per-word fractional costs are
$7/2,11/2,5/2,3,1$. Exhausting $\binom{42}{2}=861$ vocabularies certifies
the integer value. These examples demonstrate a relaxation gap; they do
not supply a universal approximation ratio for rounding.

\subsection{Application to permitted paths on actual higher-genus dessins}
Set $m=2g+1$ for $g=2,4$, and retain the existing clean dessin of
$y^2=x^m-1$, $\beta=x^m$. On $H=\{0,\ldots,2m-1\}$ put
\[
 \sigma=(0\ \cdots\ m-1)(m\ \cdots\ 2m-1),\qquad
 \alpha(h)=h+m\pmod{2m},\qquad \tau=\sigma^{-1},\quad\phi=\alpha\tau.
\]
There are two vertex cycles, $m$ edges and one face, so
$2-2g=2-m+1$. Its polynomial algebra is the infinite locally gentle
algebra
\[
 A_g=kQ/(a_h a_{\phi h}:h\in H),\qquad
 a_h:[h]\longrightarrow[\tau h].
\]
Here products are written in travel order: $ab$ traverses $a$ first and
then $b$. A left-operator matrix convention reverses these products.
Define the nonzero permitted word
$P(h,r)=a_h a_{\tau h}\cdots a_{\tau^{r-1}h}$. Our finite multiset is
\[
 \big(P(0,m+2),P(0,m+2),P(m-1,4),
      P(m,m+2),P(2m-1,4)\big).
\]
Each arrow label is one alphabet symbol, even when its printed label has
several characters. The duplicate first word retains its multiplicity.
We allow every occurring subpath of length two or three as an extra
candidate. The LP bound consequently concerns this explicitly restricted
candidate set, not arbitrary-length vocabularies.
\begin{center}\small
\begin{tabular}{rrrrrr}
\toprule
$g$ & $m$ & $|C|$ & symbol count & $L_2=I_2$ & subsets checked\\
\midrule
2&5&20&29&19&190\\
4&9&36&41&31&630\\
\bottomrule
\end{tabular}
\end{center}
The exact witnesses attain these values, and Det/Bias/Int all attain
them on the saved LP solutions. A token names a finite permitted subpath;
concatenation recovers every original arrow, hence its quiver endpoints
and original path length. This is finite corpus compression, not a
finite-dimensional quotient of $A_g$. Arbitrary token concatenations may
fail composability and need their own endpoint checks. Neither an algebra
homomorphism on a freely generated token algebra nor a representation of
arbitrary completed formal sums is asserted. The completed order still
has rank $2m^2$ over $k[[u]]$; its rank, coefficients, multiplication,
relations and modules are not reduced by counting tokens.

The companion laboratory at
\url{http://127.0.0.1:8765/convextok.html} exposes the position DAG,
literal fractional flows, vocabularies and discrete segmentations.
Its five corpora and sixteen budget cases have sixteen exact
primal/dual certificates and 2,312 exhausted integer vocabularies in
total. The browser independently recomputes segmentation costs from
local arrays; it does not solve a new LP when a user changes a vocabulary.
The primary paper's NLP training and large GPU experiments are not
reproduced. Compression of a finite path corpus proves no module or
Fukaya equivalence, arithmetic determinant comparison, or GRH statement.
